Dielektrické ztráty jsou přeměnou elektrické energie na teplo. Podle příčiny se dělí na:
\begin{itemize}
\item \textbf{Vodivostní:} pohyb volných nosičů náboje v objemu nebo po povrchu izolantu; vznikají při stejnosměrném i střídavém poli.
\item \textbf{Polarizační:} zpoždění relaxační polarizace za změnami pole; energie se spotřebovává při přemisťování iontů nebo natáčení dipólů.
\item \textbf{Ionizační:} ionizace plynu a výboje, např. v plynových dutinách izolantu při dostatečně silném poli.
\end{itemize}
